\section*{Capítulo \ref{ch:g2:pushes-and-pulls} --- Empujones y tirones: las fuerzas}

\begin{solution}{exo:g2:pushes-and-pulls:1}
Empujón; tirón; empujón; tirón; empujón (apretar es empujar por los
dos lados).
\end{solution}

\begin{solution}{exo:g2:pushes-and-pulls:2}
Dónde agarra la \omterm{def:g2:pushes-and-pulls:force}{fuerza} el objeto, hacia dónde actúa y cuánta
intensidad tiene (por la longitud de la flecha).
\end{solution}

\begin{solution}{exo:g2:pushes-and-pulls:3}
Poner en movimiento, parar, desviar y cambiar la forma. Los ejemplos
varían --- por ejemplo: dar un capirotazo a una canica (poner en
movimiento), atrapar un disco volador (parar), hacer rebotar una
pelota en la pared (desviar), aplastar una bola de nieve (forma).
\end{solution}

\begin{solution}{exo:g2:pushes-and-pulls:4}
Pararla: empujar en contra de su movimiento, con el dedo hacia la
canica. Acelerarla: empujar por detrás, en el sentido en que ya rueda.
Desviarla: empujar de lado, cruzando su camino.
\end{solution}

\begin{solution}{exo:g2:pushes-and-pulls:5}
Recuperan su forma: la esponja y la goma elástica. Se quedan con la
marca: la plastilina y la nieve blanda.
\end{solution}

\begin{solution}{exo:g2:pushes-and-pulls:6}
El empujón suave: una flecha corta apuntando a la derecha. El empujón
fuerte: una flecha larga apuntando a la izquierda. Longitudes
distintas, sentidos contrarios.
\end{solution}

\begin{solution}{exo:g2:pushes-and-pulls:7}
El tirón de la Tierra. No necesita tocar --- tira de la cuchara a
través del \omterm{def:g2:air-around-us:air}{aire}, y siempre hacia abajo.
\end{solution}

\begin{solution}{exo:g2:pushes-and-pulls:8}
El tirón de la Tierra (un calcetín que se cae al suelo) y el tirón del
\omterm{def:g1:magnets:magnet}{imán} (el \omterm{def:g1:magnets:magnet}{imán} del frigorífico sujeta un dibujo a través del papel).
\end{solution}

\begin{solution}{exo:g2:pushes-and-pulls:9}
Dos flechas sobre la soga, igual de largas y apuntando en sentidos
contrarios. La cinta se queda quieta cuando los dos tirones son
\emph{igual de intensos} --- se anulan el uno al otro. En el momento
en que un equipo tira más fuerte, gana su flecha y la cinta se mueve
hacia ese lado.
\end{solution}

\begin{solution}{exo:g2:pushes-and-pulls:10}
La misma \omterm{def:g2:pushes-and-pulls:force}{fuerza} en los tres instantes: el tirón de la Tierra, con la
flecha apuntando recta hacia abajo --- durante la subida (frenando la
pelota), arriba del todo (a punto de devolverla) y durante la bajada
(acelerándola). No deja de tirar ni un solo instante.
\end{solution}
